\ifx\pdfminorversion\undefined\else\pdfminorversion=4\fi
\documentclass[aspectratio=43,t]{beamer}
%\documentclass[aspectratio=43,t,handout]{beamer}

% English version FAU Logo
\usepackage[english]{babel}
% German version FAU Logo
%\usepackage[ngerman]{babel}

\usepackage[utf8]{inputenc}
\usepackage[T1]{fontenc}
\usepackage{amsmath,amssymb}
\usepackage{graphicx}
\usepackage{listings}
\usepackage[backend=biber,sorting=none,doi=true,style=ieee]{biblatex}

% Options:
%  - inst:      Institute
%                 med:      MedFak FAU theme
%                 nat:      NatFak FAU theme
%                 phil:     PhilFak FAU theme
%                 rw:       RWFak FAU theme
%                 rw-jura:  RWFak FB Jura FAU theme
%                 rw-wiso:  RWFak FB WISO FAU theme
%                 tf:       TechFak FAU theme
%                 tf-hscd:  TechFak Co-Design FAU theme
%  - image:     Cover image on title page
%  - plain:     Plain title page
%  - longtitle: Title page layout for long title
\usetheme[%
  inst=tf-hscd,%
  longtitle%
]{fau}

% Enable semi-transparent animation preview
\setbeamercovered{transparent}


\lstset{%
  language=C++,
  tabsize=2,
  basicstyle=\tt,
  keywordstyle=\color{blue},
  commentstyle=\color{green!50!black},
  stringstyle=\color{red},
  numbers=left,
  numbersep=0.5em,
  xleftmargin=1em,
  numberstyle=\tt
}


\defbibheading{bibliography}{}
\addbibresource[label=primary]{references.bib}
\nocite{*}


% Title, authors, and date
\title[Short Title]{Title}
\subtitle{Multi-Core Architectures and Programming}
\author[Oliver Reiche]{Oliver Reiche}
% English version
\institute[Hardware/Software Co-Design]{Hardware/Software Co-Design, Friedrich-Alexander University Erlangen-Nürnberg}
% German version
%\institute[Hardware-Software-Co-Design]{Hardware-Software-Co-Design, Friedrich-Alexander-Universität Erlangen-Nürnberg}
\date{\today}
% Set additional logo (overwrites FAU seal)
%\logo{\includegraphics[width=.15\textwidth]{themefau/art/tf/hscd/codesign.pdf}}


\begin{document}
  % Title
  \maketitle

  { % Motivation
    \setbeamertemplate{footline}{}
    \begin{frame}[noframenumbering]{Motivation}
      Some motivational example\dots
    \end{frame}
  }

  { % Outline
    \setbeamertemplate{footline}{}
    \begin{frame}[noframenumbering]{Outline}
      \tableofcontents
    \end{frame}
  }

  % Body
  \section{Lists}
  \begin{frame}{Lists}
    \begin{block}{List of Items}
      \begin{itemize}
        \item one
        \item two
        \item three
        \item \dots
      \end{itemize}
    \end{block}

    \pause

    \begin{block}{Numbered List of Items}
      \begin{enumerate}
        \item<2-> one
        \item<3-> two
        \item<4-> three
        \item<5-> \dots
      \end{enumerate}
    \end{block}
  \end{frame}

  \section{Listings}
  \begin{frame}[fragile]{Listings}
    \begin{lstlisting}[language=C++]
#include <iostream>

int main() {
  std::cout << "Hello World!" << std::endl;
  return 0;
}
\end{lstlisting}
\end{frame}

  { % Questions?
    \setbeamertemplate{footline}{}
    \begin{frame}[c,noframenumbering]
      \begin{center}
        Thanks for listening.\\
        {\bf Any questions?}
      \end{center}
    \end{frame}

    % References
    \section*{References}
    \begin{frame}[allowframebreaks,noframenumbering]{References}
      \printbibliography
    \end{frame}
  }
\end{document}

